\section{Candidate-map admissibility}
\label{sec:candidates}

A numerical workflow may construct several approximations over one total interval. Complete positivity belongs to the particular linear map that is advanced. We therefore separate candidate construction from error estimation and from the nominal RK formula.

Write $\Phi_\theta(a,c)$ for the phase-covariant map in Eq.~\eqref{eq:direct_map}; the maps below share the same stationary population parameter $\theta$.

\begin{proposition}[Composition calculus]
\label{prop:composition_calculus}
For two phase-covariant maps with the same $\theta$,
\begin{equation}
 \Phi_\theta(a_2,c_2)\circ\Phi_\theta(a_1,c_1)
 =\Phi_\theta(a_2a_1,c_2c_1).
 \label{eq:composition_calculus}
\end{equation}
Consequently, $n$ equal RK substeps over total coordinates $(x,u,\nu)$ have multipliers
\begin{equation}
 a_n=R(-x/n)^n,
 \qquad
 c_n=R[-(x/2+u-\ii\nu)/n]^n.
 \label{eq:n_substep_multipliers}
\end{equation}
\end{proposition}

\begin{proof}
Substitution into Eq.~\eqref{eq:direct_map} shows that the population deviation from the common fixed point is multiplied by $a_2a_1$, while the coherence is multiplied by $c_2c_1$.
\end{proof}

Let $\mathcal A_R$ be the full one-step CPTP set on a fixed one-way generator ray. If the substep population multiplier $a_0=R(-x/n)$ is nonnegative, then $a_0^n\le1$ exactly when $a_0\le1$, and $a_0^n-|c_0|^{2n}$ has the same sign as $a_0-|c_0|^2$. Hence
\begin{equation}
 \Phi_R(x/n)^n\ \text{is CPTP}
 \quad\Longleftrightarrow\quad
 \Phi_R(x/n)\ \text{is CPTP}.
 \label{eq:no_repair_substep}
\end{equation}
Classical RK4 satisfies $R_4(-x)>0$ for every $x\ge0$ because
\begin{equation}
 24R_4(-x)=(x^2-2x)^2+8(x^2-3x+3)>0.
\end{equation}
Hence the total-step admissible set of an $n$-substep RK4 candidate is exactly
\begin{equation}
 \mathcal A_4^{(n)}=n\mathcal A_4.
 \label{eq:dilated_candidate_set}
\end{equation}

\begin{proposition}[Substep-ladder overlap of positive components]
\label{prop:substep_ladder}
Suppose a nontrivial positive-step component is $\mathcal I=[x_-,x_+]$ with $0<x_-<x_+$. Then
\begin{equation}
 \mathcal I^{(n)}\cap\mathcal I^{(n+1)}\ne\varnothing
 \quad\Longleftrightarrow\quad
 \frac{x_+}{x_-}\ge\frac{n+1}{n}.
 \label{eq:ladder_overlap}
\end{equation}
For the full-step/two-half-step pair ($n=1$), the nontrivial positive-step components are disjoint exactly when $x_+/x_-<2$. The same factor-two specialization is equivalent to the existence of initial proposals for which repeated halving skips $\mathcal I$. More generally, a fixed reduction factor $\rho$ hits the component for every proposal above it if and only if $\rho\ge x_-/x_+$.
\end{proposition}

\begin{proof}
Equation~\eqref{eq:dilated_candidate_set} gives the corresponding positive components $\mathcal I^{(n)}=[nx_-,nx_+]$ and $\mathcal I^{(n+1)}=[(n+1)x_-,(n+1)x_+]$. Their overlap is equivalent to $(n+1)x_-\le nx_+$. The reduction statement follows by comparing two consecutive proposals $\rho^jx_0>x_+$ and $\rho^{j+1}x_0<x_-$.
\end{proof}

At $\varpi=2$,
\begin{equation}
 \mathcal A_4=\{0\}\cup[0.7447824278959\ldots,1.2703346269739\ldots],
 \label{eq:direct_varpi2_interval}
\end{equation}
whereas
\begin{equation}
 \mathcal A_4^{(2)}=\{0\}\cup[1.4895648557918\ldots,2.5406692539478\ldots].
 \label{eq:fine_varpi2_interval}
\end{equation}
The sets share only the zero step; their nontrivial positive-step components are disjoint. Thus $H=1$ gives a CPTP full candidate and a non-CPTP fine candidate, whereas $H=2$ reverses those statuses. The step-doubling pair is also the most demanding consecutive pair: at $\varpi=2$, the corresponding positive components $\mathcal I_4^{(n)}$ and $\mathcal I_4^{(n+1)}$ overlap for every $n\ge2$. Positive-component disjointness and the halving-skip condition end together at
\begin{equation}
 \varpi_{1/2}=2.248254604514980\ldots,
 \label{eq:halving_threshold}
\end{equation}
whose exact algebraic certificate is given in the Supplementary Material.

\subsection{Affine candidates and the extrapolation trap}

The CPTP set is convex. If $\Phi_j$ are CPTP, $w_j\ge0$, and $\sum_jw_j=1$, then $\sum_jw_j\Phi_j$ is CPTP. Trace preservation survives every affine combination with weights summing to one, but negative weights need not preserve complete positivity.

For classical RK4, let $\Phi_{\rm c}(x)$ be the coarse full-step map and $\Phi_{\rm f}(x)=\Phi_{\rm RK4}(x/2)^2$ the fine map. The standard local Richardson extrapolate built from fourth-order RK4 is \cite{Szyszkowicz1985Richardson,ZlatevEtAl2017Richardson,HavasiKazemi2018Richardson}
\begin{equation}
 \Phi_{\rm ext}(x)=\frac{16\Phi_{\rm f}(x)-\Phi_{\rm c}(x)}{15}.
 \label{eq:richardson_map}
\end{equation}
The affine candidate is itself generated by the stability polynomial
\begin{equation}
 R_{\rm ext}(z)=\frac{16R_4(z/2)^2-R_4(z)}{15}
 =\sum_{j=0}^{5}\frac{z^j}{j!}+\frac{z^6}{864}+\frac{z^7}{8640}+\frac{z^8}{138240}.
 \label{eq:richardson_stability}
\end{equation}
Its first exponential defect is
\begin{equation}
 R_{\rm ext}(z)-\e^z=-\frac{z^6}{4320}+O(z^7).
 \label{eq:richardson_defect}
\end{equation}
Thus the affine construction cancels RK4's fifth-order exponential defect and has linear order five. On the one-way zero-dephasing boundary the exact margin vanishes identically, so the same cancellation removes the constituents' leading inward $x^5$ margin and exposes the next defect order. The signed first-defect theorem therefore gives the local physicality failure without a separate Choi expansion:
\begin{equation}
 M_{\rm ext}(x;\varpi)
 =-\frac{G_6(\varpi)}{4320}x^6+O(x^7),
 \qquad
 G_6=\frac{q(q^2-18q+48)}{32},
 \label{eq:richardson_frequency}
\end{equation}
where $q=1+4\varpi^2$. Its two positive orientation thresholds are
\begin{equation}
 \varpi_{\mathrm{ext},-}=\sqrt{2-\frac{\sqrt{33}}4},
 \qquad
 \varpi_{\mathrm{ext},+}=\sqrt{2+\frac{\sqrt{33}}4}.
 \label{eq:richardson_frequency_thresholds}
\end{equation}
The extrapolate is locally outward for $0\le\varpi<\varpi_{\mathrm{ext},-}$, inward between the thresholds, and outward again above $\varpi_{\mathrm{ext},+}$. The Dormand--Prince principal formula has the opposite sequence because it has the same $G_6$ but the opposite sign of $\eta_6$. This all-frequency laboratory-frame orientation remains representation dependent, as does every nonzero-$\varpi$ direct-RK statement. By contrast, the global theorem below is proved at $\varpi=0$, where no rotating-frame change is available or needed; its extrapolation failure is therefore frame independent. Supplemental Fig.~S1 displays the three frequency-orientation bands together.

\begin{theorem}[CPTP loss under Richardson extrapolation of RK4]
\label{thm:richardson}
On nonrotating one-way damping, both $\Phi_{\rm c}(x)$ and $\Phi_{\rm f}(x)$ are CPTP for $0\le x\le\alpha_4$. For every $0<x\le\alpha_4$, the extrapolated map \eqref{eq:richardson_map} is non-CPTP. Its margin factors as
\begin{equation}
 M_{\rm ext}(x)
 =-\frac{x^6P_{10}(x)}{1252412463513600}
 =-\frac{31}{138240}x^6+O(x^7),
 \label{eq:richardson_margin}
\end{equation}
where $P_{10}$ is the degree-ten integer polynomial listed in the Supplementary Material and is strictly positive on $[0,\alpha_4]$. Globally, the extrapolated map later re-enters the CPTP set on the remote interval
\begin{equation}
 [6.034523958649\ldots,6.459127767826\ldots].
 \label{eq:richardson_remote}
\end{equation}
\end{theorem}

\begin{proof}
Equation~\eqref{eq:richardson_defect} and Theorem~\ref{thm:defect} give the leading coefficient in Eq.~\eqref{eq:richardson_margin}. Exact substitution gives the factor $P_{10}$. A Sturm sequence isolates its first positive zero at $6.034523958649\ldots>\alpha_4$, proving outward orientation over the common CPTP interval of the constituents. The complete constraint sign chart then isolates the remote component \eqref{eq:richardson_remote} and its population-ceiling exit.
\end{proof}

The theorem applies when the extrapolated candidate is actually advanced. Many step-doubling codes advance the fine map and use the difference only for error estimation. The general conclusion is candidate-specific: neither the nominal RK method nor the physicality of constituent candidates determines the physicality of a negative-weight affine output.

\begin{figure*}[t]
 \centering
 \makebox[\textwidth][c]{\includegraphics[width=\JCPFigureWidth]{Fig3_candidate_maps.pdf}}
 \caption{Candidate-map complete positivity. (a) At $\varpi=2$, the nontrivial positive-step CPTP components of the full RK4 and two-half-step candidates are disjoint; both sets also contain the zero step. The marked total steps reverse which candidate is CPTP. (b) On nonrotating one-way damping, the coarse and fine candidates are CPTP through $x=\alpha_4$, whereas the negative-weight Richardson extrapolate has a strictly negative Choi margin throughout that common interval.}
 \label{fig:candidate_maps}
\end{figure*}
